\chapter{超块抽屉：同一格不同标签}
\label{ch:superblock}

\section{抽屉比喻}

\begin{metaphor}[共享收纳盒]
物理上多出来的 padding 被挖成几个小格。\\
格子编号相同，但 Hom/Chain/Native 往里放的纸条\textbf{含义不同} —— 读之前必须先看房门 flag。
\end{metaphor}

\begin{figure}[htbp]
\centering
\begin{tikzpicture}[font=\footnotesize]
  \draw[thick] (0,0) rectangle (11,3.2);
  \node[anchor=west, font=\bfseries] at (0.2,2.9) {ext\_padding 挖出的 topo 区};
  \node[softbox=gray, minimum width=2.4cm] at (1.5,1.6) {seed\\u32};
  \node[softbox=gray, minimum width=1.6cm] at (4.2,1.6) {variant\\u8};
  \node[softbox=gray, minimum width=2.6cm] at (7.0,1.6) {reserved\\u16+};
  \node[font=\tiny, gray] at (5.5,0.4)
    {seed=骰子/锁芯/炼金种子；u16=calib 或 L 或 stride…};
\end{tikzpicture}
\caption{盲读 reserved $=$ 把药说明书当菜谱。}
\label{fig:drawer}
\end{figure}

\section{逐代纸条内容}

\begin{center}
\small
\begin{tabular}{@{}llll@{}}
\toprule
格 & Hom/TopoPH & Chain & Native \\
\midrule
seed & Gear 骰子 & LFSR 初值 & Embed 种子 \\
variant & 0/1 或 3/4 & 2 & 5/6 \\
u16 & calib\_permille & stride\_log $L$ & event\_stride \\
其后 & shift 等 & quant、beta1 位 & $k$、persist 位 \\
\bottomrule
\end{tabular}
\end{center}

\section{开仓时写入 vs 备份时读回}

InitRepo：掷种子、标定、写纸条。\\
RunBackup：必须原样读回 —— 否则等于半夜换刀。

\begin{pitfall}[误清 FastCDC]
Topo 仓恢复时若把 flag 误清成 Fast，后果是灾难级贴纸错乱。正式代码有护栏。
\end{pitfall}
